%% kapa: exam=1 round=31 year=2020 period=1 form=A subject=civil_law subject_ko=민법 seq=33 number=33
%% source: Q-Net 감정평가사 기출문제 (cst00309) artlSeq=5207219 file=extracted/1008/5207219/A형/A형/2020년도 제31회 감정평가사 1차 시험_1교시_A형.pdf
%% conversion: kordoc 4.24.51 -> tools/gen1.py (auto); review=proofread-sonnet; flags=none
\begin{question}{33}
첨부에 관한 설명으로 옳지 않은 것은? (다툼이 있으면 판례에 따름)
\begin{choices}
  \item 타인이 그의 권원에 의하여 부동산에 부속한 물건은 이를 분리하여도 경제적 가치가 없으면 부동산소유자의 소유로 한다.
  \item 저당권의 효력은 다른 사정이 없으면 저당부동산에 부합된 물건에 미친다.
  \item 동일인 소유의 여러 동산들이 결합하는 것은 부합이 아니다.
  \item 부합의 원인은 인위적이든 자연적이든 불문한다.
  \item 타인의 동산에 가공한 때에는 가공물의 소유권은 가공자의 소유로 한다.
\end{choices}
\end{question}
\answer{5}
